\documentclass[10pt,a4paper]{amsart}
\usepackage[margin=19mm]{geometry}
\usepackage{amsmath,amssymb,mathtools}
\usepackage[scr=rsfs,cal=euler]{mathalfa}
\usepackage{xfrac}
\usepackage[unicode,hidelinks,bookmarks=false]{hyperref}
\newcommand{\ie}{i.e.\ }
\newcommand{\cf}{cf.\ }
\newcommand{\resp}{respectively\ }
\newcommand{\Subsection}{\subsection}
\providecommand{\nobreakdash}{\nobreak\hbox{-}}
\setlength{\emergencystretch}{2em}
\multlinegap0pt
\usepackage{amssymb}
\usepackage{mathtools}
\newtheorem{proposition}{Proposition}
\def\H{\mathcal{H}}
\title{Morse diagrams, Murasugi sums, and the mapping class group}
\author{Jack Brand, David Gay,, Joan Licata}
\date{Source: arXiv:2604.00448v1 (CC BY 4.0)}
\begin{document}
\maketitle

\noindent\emph{Source and licence.} Morse diagrams, Murasugi sums, and the mapping class group. Version arXiv:2604.00448v1 (CC BY 4.0), distributed by arXiv under the Creative Commons Attribution 4.0 International licence, http://creativecommons.org/licenses/by/4.0/, as recorded in the arXiv licence metadata for this exact version and shown on the abstract page https://arxiv.org/abs/2604.00448v1. Full licence notice: \texttt{licenses/arxiv-2604.00448-CC-BY-4.0.txt}.\par\smallskip

\noindent\textit{Editorial scope.} Counted unit: the article's proposition prop:MakeNice (source0.tex lines 415--417). The article's complete original proof of it is delivered with it (source0.tex lines 498--506, one \texttt{proof} environment of 9 lines, which the article places away from the statement). No mathematics is written, repaired or re-derived: the statement and the proof are the article's own lines, in the article's own notation, and no retained line is substituted or re-flowed.  No further source-local context is needed: every cross-reference of every retained line resolves inside the two delivered spans.  Nothing was omitted from any retained span, so the package carries no omission record.  No retained line contains a citation, so the wrapper carries no bibliography and no bibliography entry.  No retained line contains a figure environment, an image-inclusion command or drawing code, so no image is delivered and the package declares no asset dependency.  Editorial formatting only, in addition: an AMS \texttt{amsart} preamble that restates the article's own declarations and macros used by the retained lines, namely the environments \texttt{proposition} on separate counters, together with the standard packages those lines need.  The retained environments print in this document as Proposition 1 in order of appearance, whereas the article prints the same items with the numbering of its own full document.  The complete native source of the article as distributed in the arXiv e-print for this exact version is bundled unchanged as \texttt{sources/197590/source0.tex}.
\begin{proposition}\label{prop:MakeNice} 
If $D$ is a Murasugi polygon on $\Sigma_i$, there exists a handle structure such that the points $p_j^i$ are starlike.   
\end{proposition}

\begin{proof}[Proof of \ref{prop:MakeNice}]
Let  $D$ denote the star-shaped graph associated to the Murasugi polygon and let $v$ denote the $n$-valent vertex.   We will build a handle structure $\H$ for $\Sigma$ such that   $D\cap\H=v$. 
    
    Enumerate the closures of the components of $\Sigma\setminus\delta$  as $\Sigma_{1},\dots,\Sigma_{k}$, allowing $k=1$ if $\Sigma\setminus\delta$ is  connected. The boundary of each $\Sigma_{i}$ has arcs that glue to $\delta$; let $q_j$ be the points on these arcs that glue to $q\in\Sigma$.
    
   For each component $\Sigma_i$, choose a handle structure.  Label the wedge point for the handle structure on $\Sigma_i$  by $v_i$.  Each $\Sigma_i$ has one or more copies of $v$ on its boundary; isotope the handle structure on each component so that $v_i$ lies on a copy of $v\in \partial \Sigma_i$.  For each unused copy of $v\in \partial \Sigma_i$, choose a properly embedded arc whose interior is disjoint from the handle structure and that connects $v$ to $v_i$.  We claim that the image of these curves and the component-wise handle structures include a complete set of cores for $\Sigma$.
      
 By construction, the points $v_i$ are all identified with $v$, forming a single wedge point for the cores.  To see that we have enough cores for a handle structure, let $C$ be a simple closed curve in $\Sigma$.  Each time $C$ crosses an edge of $D$, isotope $C$ so that it crosses at $v$.  This decomposes $C$ into a collection of arcs, each of which is either a closed curve or a properly embedded curve in some $\Sigma_i$.  Each closed component deformation retracts onto the cores of the corresponding handle structure, while each curve with distinct endpoints in $\Sigma_i$ is isotopic to a concatenation of the added curves.  Finally, co-cores can be added to produce the desired handle structure.  
 \end{proof}

\end{document}
